\section{EQUIVALENCE RELATIONS AND QUOTIENT SETS}

\begin{enumerate}
\def\labelenumi{\arabic{enumi}.}
\tightlist
\item
  Let \((\mathbf{A}_t)_{i \in \mathbf{I}}\) be a partition of a set \(\mathbf{E}\) . The relation \(\mathbf{R}\{x, y\} :\) ``there exists \(i \in \mathbf{I}\) such that \(x \in \mathbf{A}_t\) and \(y \in \mathbf{A}_t\)'' between two generic elements \(x, y\) of \(\mathbf{E}\) satisfies the following conditions:
\end{enumerate}

\begin{enumerate}
\def\labelenumi{(\alph{enumi})}
\item
  \(\mathbf{R}\{x,x\}\) is an identity (reflexivity of \(\mathbf{R}\) ).
\item
  \(\mathbf{R}\{x,y\}\) and \(\mathbf{R}\{y,x\}\) are equivalent (symmetry of \(\mathbf{R}\) ).
\item
  The relation ``R\{x,y\} and R\{y,z\}'' implies R\{x,z\} (transitivity of R).
\end{enumerate}

If C denotes the subset of \(E \times E\) defined by the relation R, the conditions (a), (b), (c) are respectively equivalent to the following conditions: \((a') \Delta \subset C; (b') \overline{C} = C; (c') C \circ C \subset C.\) From \((a')\) and \((c')\) it follows that \(C \circ C = C.\)

\begin{enumerate}
\def\labelenumi{\arabic{enumi}.}
\setcounter{enumi}{1}
\tightlist
\item
  Conversely, let \(R\{x,y\}\) be a reflexive, symmetric, and transitive relation, and let C be its graph in \(E \times E\) . Then the image F of E under the mapping \(x \to C(x)\) of E into \(\mathfrak{P}(E)\) is a partition of E, and the relation ``there exists a subset \(X \in F\) such that \(x \in X\) and \(y \in X\) '' is equivalent to \(R\{x, y\}\) .
\end{enumerate}

Every relation which satisfies conditions (a), (b), and (c) is called an equivalence relation on E. The partition \(\mathfrak{F}\) which it defines, considered as a subset of \(\mathfrak{P}(\mathbf{E})\) , is called the quotient set of E by the relation R, and is denoted by E/R; its elements are called equivalence classes with respect to R. The mapping \(x \to \mathrm{C}(x)\) of E onto E/R, which maps each element \(x\) of E to the equivalence class which contains \(x\) , is called the canonical mapping of E onto E/R.

The relation of equality x = y is an equivalence relation. The canonical mapping of E onto the corresponding quotient set is just \(x \to \{x\}\) , and is bijective.

If R is an equivalence relation, the notation `` \(x \equiv y \pmod{R}\) '' is sometimes used as a synonym for R \(\{x, y\}\) ; it is read `` \(x\) is equivalent to \(y\) modulo R''.

\begin{enumerate}
\def\labelenumi{\arabic{enumi}.}
\setcounter{enumi}{2}
\tightlist
\item
  On a product set \(E \times F\) , the relation `` \(pr_{1}z = pr_{1}z'\) '' between z and \(z'\) is an equivalence relation R, and the quotient set \((\mathbf{E} \times \mathbf{F}) / \mathbf{R}\) can be put in one-to-one correspondence with E (this is the origin of the name quotient set).
\end{enumerate}

More generally, let \(f\) be a mapping of a set \(\mathbf{E}\) into a set \(\mathbf{F}\) . Then the relation '' \(f(x) = f(y)\) '' is an equivalence relation on \(\mathbf{E}\) . If we denote this relation by \(\mathbf{R}\) , the mapping \(z \to \overline{f}^1(z)\) (where \(\overline{f}^1(z)\) is considered as an element of \(\mathbf{E}/\mathbf{R}\) ) is a bijection of \(f(\mathbf{E})\) onto \(\mathbf{E}/\mathbf{R}\) .

It follows that \(f\) may be considered as the composition of the following three mappings, in the given order:

\begin{enumerate}
\def\labelenumi{(\arabic{enumi})}
\item
  the canonical mapping of the subset \(f(\mathbf{E})\) of \(\mathbf{F}\) into the set \(\mathbf{F}\) ;
\item
  the bijective mapping of E/R onto \(f(\mathbf{E})\) , whose inverse has just been defined;
\item
  the canonical mapping of E onto E/R.
\end{enumerate}

This decomposition of a mapping is called its canonical decomposition or canonical factorization.

\begin{enumerate}
\def\labelenumi{\arabic{enumi}.}
\setcounter{enumi}{3}
\item
  Every equivalence relation R on a set E may be defined by means of a mapping as in the previous subsection; for, if C is the graph of R, the relation ``C(x) = C(y)'' is equivalent to R\{x, y\}.
\item
  Let R be an equivalence relation on a set E, and let A be a subset of E. Then the relation \(R\{x, y\}\) between two generic elements x, y of A is an equivalence relation on A, called the relation induced by R on A, and denoted by \(R_{A}\) . Let f be the canonical mapping of E onto E/R and let g be that of A onto \(A/R_{A}\) . By making an element of E/R and an element of \(A/R_{A}\) correspond if they are the images of the same element of E under f and g respectively, we have a one-to-one correspondence between the image \(f(A)\) of A under f and the quotient set \(A/R_{A}\) . If \(\varphi\) denotes the canonical mapping of A into E, this correspondence is realized by the mapping \(z \to f(\varphi(\overline{g}^1(z)))\) and its inverse, both of which are also called canonical.
\item
  A subset \(\mathbf{A}\) of \(\mathbf{E}\) is said to be saturated with respect to the equivalence relation \(\mathbf{R}\) if for each \(x \in \mathbf{A}\) the equivalence class of \(x\) with respect to \(\mathbf{R}\) is contained in \(\mathbf{A}\) . In other words, the saturated sets with respect to \(\mathbf{R}\) are unions of equivalence classes with respect to \(\mathbf{R}\) . If \(f\) is the canonical mapping of \(\mathbf{E}\) onto \(\mathbf{E} / \mathbf{R}\) , then a set is saturated if it is of the form \(\overline{f}^1(\mathbf{X})\) , where \(\mathbf{X} \subset \mathbf{E} / \mathbf{R}\) .
\end{enumerate}

Let A be a subset of E. The intersection of the saturated sets which contain A is \(\overline{f}^{1}(f(A))\) . This set may also be defined as the union of the equivalence classes of the elements of A, and is called the saturation of A (with respect to R).

\begin{enumerate}
\def\labelenumi{\arabic{enumi}.}
\setcounter{enumi}{6}
\tightlist
\item
  Let \(\mathrm{P}\{x, y, z\}\) be a relation in which there appears a generic element \(x\) of \(\mathbf{E}\) . Then \(\mathrm{P}\) is said to be compatible (in \(x\) ) with the equivalence relation \(\mathbf{R}\) if the relation `` \(\mathrm{P}\{x, y, z\}\) and \(x \equiv x' (\text{mod } \mathbf{R})\) '' implies \(\mathrm{P}\{x', y, z\}\) .
\end{enumerate}

Let \(f\) be the canonical mapping of \(\mathbf{E}\) onto \(\mathbf{E} / \mathbf{R}\) , and let \(t\) be a generic element of \(\mathbf{E} / \mathbf{R}\) . The relation ``there exists \(x \in \overline{f}^1(t)\) such that \(\mathrm{P}\{x, y, z\}\)'' is then equivalent to ``for all \(x \in \overline{f}^1(t)\) , \(\mathrm{P}\{x, y, z\}\) ; the latter is a relation between \(t, y, z\) , said to be induced by \(\mathrm{P}\) on passing to the quotient (with respect to \(x\) ). If we denote it by \(\mathrm{P}'\{t, y, z\}\) , then \(\mathrm{P}\{x, y, z\}\) is equivalent to \(\mathrm{P}'\{f(x), y, z\}\) .

There are analogous definitions for a relation involving any number of arguments, and for the case in which the relation is compatible with R in several of its arguments. For example, if A is a subset of E, to say that the relation `` \(x \in A\) '' is compatible (in \(x\) ) with R is equivalent to saying that A is saturated with respect to R. If \(\varphi\) is a mapping of E into a set F, to say that the functional relation `` \(y = \varphi(x)\) '' is compatible (in \(x\) ) with R is to say that \(\varphi\) is constant on each equivalence class with respect to R. Passing to the quotient, R therefore induces a relation between \(y\) and a generic element \(t\) of E/R; this relation is functional in \(y\) and so determines a mapping \(\varphi'\) of E/R into F, satisfying the identity \(\varphi(x) = \varphi'(f(x))\) .

\begin{enumerate}
\def\labelenumi{\arabic{enumi}.}
\setcounter{enumi}{7}
\item
  Let R be an equivalence relation on a set E, let S be an equivalence relation on a set F, and let f be a mapping of E into F. The mapping f is said to be compatible with R and S if the relation \(x \equiv x' \pmod{R}\) implies \(f(x) \equiv f(x') \pmod{S}\) . If g is the canonical mapping of F onto F/S, then the composite function \(g \circ f\) has the same value at all elements of an equivalence class z with respect to R; if we denote this common value by \(h(z)\) , then h is a mapping of E/R into F/S, and is said to be induced by f on passing to the quotients.
\item
  Let R be an equivalence relation on E, and let S be an equivalence relation on E/R. If f is the canonical mapping of E onto E/R, then `` \(f(x) \equiv f(y) \pmod{S}\) '' is an equivalence relation T on E. An equivalence class with respect to T is therefore the union in E of equivalence classes with respect to R which are equivalent to each other with respect to S; and the relation `` \(x \equiv y \pmod{R}\) '' implies `` \(x \equiv y \pmod{T}\) ''. If \(g\) and \(\varphi\) are the canonical mappings of E/R onto (E/R)/S and E onto E/T, respectively, we obtain a one-to-one correspondence (called canonical) between (E/R)/S and E/T by making an element of (E/R)/S and an element of E/T correspond to each other if they are the images of the same element of E under \(g \circ f\) and \(\varphi\) , respectively.
\end{enumerate}

Conversely, let R and T be two equivalence relations on E such that `` \(x \equiv y \pmod{R}\) '' implies `` \(x \equiv y \pmod{T}\) ''. Then T is compatible (in the sense of no. 7) with R, both in \(x\) and in \(y\) ; by passing to the quotient E/R (with respect to \(x\) and \(y\) ), T induces an equivalence relation S on E/R. If \(f\) again denotes the canonical mapping of E onto E/R, the relation `` \(f(x) = f(y) \pmod{S}\) '' is equivalent to `` \(x \equiv y \pmod{T}\) ''. The equivalence relation S is called the quotient of T by R, and is denoted by T/R. From the previous paragraph we see that there exists a one-to-one correspondence (the canonical correspondence) between (E/R)/(T/R) and E/T.

\begin{enumerate}
\def\labelenumi{\arabic{enumi}.}
\setcounter{enumi}{9}
\tightlist
\item
  Now let E, F be any two sets, which may or may not be distinct. Let R\{x, y\} be an equivalence relation on E and let S\{z, t\} be an equivalence relation on F. Then the relation ``R\{x, y\} and S\{z, t\}'' between elements (x, z) and (y, t) of the product set E × F is an equivalence relation on E × F, called the product of R by S, and denoted by R × S. Every equivalence class with respect to R × S is the product of an equivalence class with respect to R and an equivalence class with respect to S. If u denotes a generic element of E/R, and v a generic element of F/S, then (u, v) → u × v is a bijection (called canonical) of (E/R) × (F/S) onto (E × F)/(R × S).
\end{enumerate}

